% GENERATED FILE. Edit canonical lessons, then run npm run generate:books.
% canonical-source-hash: fnv1a64:c516e2fd9e34f7cf
% canonical-lessons: MW-C216-nikalno, MW-C216-sathi, MW-C216-kos, MW-C216-musaphar, MW-C216-dero

\chapter{Setting Out}
\label{ch:mw-setting-out}
\hlchaptermodality{216}

\begin{chapteropening}
\textbf{By the end of this chapter:} \emph{I can say when I will set out and how far I will go.}

Complete the last lesson of chapter 216: I can say when I will set out and how far I will go.
\end{chapteropening}

\section[\texorpdfstring{\mw{निकळणो} (nikaḷṇo)}{निकळणो (nikaḷṇo)}]{\mw{निकळणो} (nikaḷṇo) — to set out; to come out}

\label{lesson:MW-C216-nikalno}



\begin{quote}
Before the new one: say the Marwadi for brave, then the Marwadi for crazy.
\end{quote}

\subsection*{You'll want to know: \mw{निकळणो}}
\textbf{\mw{निकळणो}} — \emph{nikaḷṇo} — \textquotedblleft{}to set out; to come out\textquotedblright{}.

\begin{quote}
\mw{हूं} \mw{काल} \mw{तड़कै} \mw{घर} \mw{सूं} \mw{निकळसूं।}
\end{quote}

\begin{quote}
\emph{hū̃ kāl taṛkai ghar sū̃ nikaḷsū̃.}
\end{quote}

\begin{quote}
Tomorrow I'll set out from home at daybreak.
\end{quote}

Past (he …): \mw{वो} \mw{तड़कै} \mw{निकळ्यो}, \emph{vo taṛkai nikaḷyo}.

\begin{cousinweb}
\textbf{\mw{निकळणो}} (\emph{nikaḷṇo}) is inherited from a reconstructed \emph{*niṣkalati}, \textquotedblleft{}comes out\textquotedblright{}; CDIAL lists the Old Marwari verb \emph{nikalaï}, with the same sense, under it (CDIAL 7478).
\end{cousinweb}

\begin{grammarlens}[title={What you've built}]
One more word for saying what comes next.
\end{grammarlens}

\subsection*{Your turn}
\begin{itemize}
\raggedright
  \item \emph{Say it:} \emph{nikaḷṇo}
  \item \emph{Say it:} \emph{nikaḷṇo}, once more
  \item \emph{Say it:} \emph{hū̃ kāl taṛkai ghar sū̃ nikaḷsū̃.} — Tomorrow I'll set out from home at daybreak.
  \item \emph{Recall:} say the Marwadi for brave, then the Marwadi for crazy, then say \emph{nikaḷṇo} again
\end{itemize}

\begin{tcolorbox}[breakable,colback=teal!4,colframe=teal!35!black,title={Before you move on}]
What does \mw{निकळणो} mean? (\textquotedblleft{}To set out; to come out\textquotedblright{}.) Say it once more.
\end{tcolorbox}
\section[\texorpdfstring{\mw{साथी} (sāthī)}{साथी (sāthī)}]{\mw{साथी} (sāthī) — a companion; a fellow traveller}

\label{lesson:MW-C216-sathi}



\begin{quote}
Before the new one: say the Marwadi for crazy, then the Marwadi for to set out.
\end{quote}

\subsection*{You'll want to know: \mw{साथी}}
\textbf{\mw{साथी}} — \emph{sāthī} — \textquotedblleft{}a companion; a fellow traveller\textquotedblright{}.

\begin{quote}
\mw{राम} \mw{म्हारो} \mw{साथी} \mw{है।}
\end{quote}

\begin{quote}
\emph{rām mhāro sāthī hai.}
\end{quote}

\begin{quote}
Ram is my travelling companion.
\end{quote}

A masculine or feminine noun, by the person meant.

\begin{cousinweb}
\textbf{\mw{साथी}} (\emph{sāthī}) is inherited from a reconstructed \emph{*sārthin}, \textquotedblleft{}having part in a company\textquotedblright{}, built on Sanskrit \emph{sārtha}, a band of travellers; CDIAL lists Marwari \emph{sāthī}, \textquotedblleft{}a comrade\textquotedblright{} (CDIAL 13366).
\end{cousinweb}

\begin{grammarlens}[title={What you've built}]
One more word for saying what comes next.
\end{grammarlens}

\subsection*{Your turn}
\begin{itemize}
\raggedright
  \item \emph{Say it:} \emph{sāthī}
  \item \emph{Say it:} \emph{sāthī}, once more
  \item \emph{Say it:} \emph{rām mhāro sāthī hai.} — Ram is my travelling companion.
  \item \emph{Recall:} say the Marwadi for crazy, then the Marwadi for to set out, then say \emph{sāthī} again
\end{itemize}

\begin{tcolorbox}[breakable,colback=teal!4,colframe=teal!35!black,title={Before you move on}]
What does \mw{साथी} mean? (\textquotedblleft{}A companion; a fellow traveller\textquotedblright{}.) Say it once more.
\end{tcolorbox}
\section[\texorpdfstring{\mw{कोस} (kos)}{कोस (kos)}]{\mw{कोस} (kos) — a kos (a distance of about two miles)}

\label{lesson:MW-C216-kos}



\begin{quote}
Before the new one: say the Marwadi for to set out, then the Marwadi for a companion.
\end{quote}

\subsection*{You'll want to know: \mw{कोस}}
\textbf{\mw{कोस}} — \emph{kos} — \textquotedblleft{}a kos (a distance of about two miles)\textquotedblright{}.

\begin{quote}
\mw{गांव} \mw{अठै} \mw{सूं} \mw{दस} \mw{कोस} \mw{है।}
\end{quote}

\begin{quote}
\emph{gā̃v aṭhai sū̃ das kos hai.}
\end{quote}

\begin{quote}
The village is ten kos from here.
\end{quote}

A masculine noun.

\begin{cousinweb}
\textbf{\mw{कोस}} (\emph{kos}) is inherited from Sanskrit \emph{krośa}, first \textquotedblleft{}a shout\textquotedblright{}, then the distance a shout carries (Monier-Williams); CDIAL lists Old Marwari \emph{kosa}, \textquotedblleft{}a distance of about two miles\textquotedblright{} (CDIAL 3611).
\end{cousinweb}

\begin{grammarlens}[title={What you've built}]
One more word for saying what comes next.
\end{grammarlens}

\subsection*{Your turn}
\begin{itemize}
\raggedright
  \item \emph{Say it:} \emph{kos}
  \item \emph{Say it:} \emph{kos}, once more
  \item \emph{Say it:} \emph{gā̃v aṭhai sū̃ das kos hai.} — The village is ten kos from here.
  \item \emph{Recall:} say the Marwadi for to set out, then the Marwadi for a companion, then say \emph{kos} again
\end{itemize}

\begin{tcolorbox}[breakable,colback=teal!4,colframe=teal!35!black,title={Before you move on}]
What does \mw{कोस} mean? (\textquotedblleft{}A kos (a distance of about two miles)\textquotedblright{}.) Say it once more.
\end{tcolorbox}
\section[\texorpdfstring{\mw{मुसाफर} (musāphar)}{मुसाफर (musāphar)}]{\mw{मुसाफर} (musāphar) — a traveller}

\label{lesson:MW-C216-musaphar}



\begin{quote}
Before the new one: say the Marwadi for a companion, then the Marwadi for a kos.
\end{quote}

\subsection*{You'll want to know: \mw{मुसाफर}}
\textbf{\mw{मुसाफर}} — \emph{musāphar} — \textquotedblleft{}a traveller\textquotedblright{}.

\begin{quote}
\mw{मुसाफर} \mw{दूर} \mw{सूं} \mw{आवै} \mw{है।}
\end{quote}

\begin{quote}
\emph{musāphar dūr sū̃ āvai hai.}
\end{quote}

\begin{quote}
The traveller comes from far away.
\end{quote}

A masculine or feminine noun, by the person meant.

\begin{cousinweb}
\textbf{\mw{मुसाफर}} (\emph{musāphar}) is usually traced to Persian \emph{musāfir}, \textquotedblleft{}a traveller\textquotedblright{}, which Persian took from Arabic; Marwadi says it with \emph{ph} and ends it in \emph{-ar}.
\end{cousinweb}

\begin{grammarlens}[title={What you've built}]
One more word for saying what comes next.
\end{grammarlens}

\subsection*{Your turn}
\begin{itemize}
\raggedright
  \item \emph{Say it:} \emph{musāphar}
  \item \emph{Say it:} \emph{musāphar}, once more
  \item \emph{Say it:} \emph{musāphar dūr sū̃ āvai hai.} — The traveller comes from far away.
  \item \emph{Recall:} say the Marwadi for a companion, then the Marwadi for a kos, then say \emph{musāphar} again
\end{itemize}

\begin{tcolorbox}[breakable,colback=teal!4,colframe=teal!35!black,title={Before you move on}]
What does \mw{मुसाफर} mean? (\textquotedblleft{}A traveller\textquotedblright{}.) Say it once more.
\end{tcolorbox}
\section[\texorpdfstring{\mw{डेरो} (ḍero)}{डेरो (ḍero)}]{\mw{डेरो} (ḍero) — a camp; a stopping place for the night}

\label{lesson:MW-C216-dero}



\begin{quote}
Before the new one: say the Marwadi for a kos, then the Marwadi for a traveller.
\end{quote}

\subsection*{You'll want to know: \mw{डेरो}}
\textbf{\mw{डेरो}} — \emph{ḍero} — \textquotedblleft{}a camp; a stopping place for the night\textquotedblright{}.

\begin{quote}
\mw{आज} \mw{रात} \mw{म्हारो} \mw{डेरो} \mw{डूंगर} \mw{कनै} \mw{है।}
\end{quote}

\begin{quote}
\emph{āj rāt mhāro ḍero ḍū̃gar kanai hai.}
\end{quote}

\begin{quote}
Tonight my camp is near the hill.
\end{quote}

A masculine noun.

\begin{cousinweb}
\textbf{\mw{डेरो}} (\emph{ḍero}) is inherited from a reconstructed \emph{*ḍēra}, \textquotedblleft{}a resting-place\textquotedblright{}; CDIAL lists Marwari \emph{ḍero}, \textquotedblleft{}a tent, a shelter\textquotedblright{} (CDIAL 5564).
\end{cousinweb}

\begin{grammarlens}[title={What you've built}]
One more word for saying what comes next.
\end{grammarlens}

\subsection*{Your turn}
\begin{itemize}
\raggedright
  \item \emph{Say it:} \emph{ḍero}
  \item \emph{Say it:} \emph{ḍero}, once more
  \item \emph{Say it:} \emph{āj rāt mhāro ḍero ḍū̃gar kanai hai.} — Tonight my camp is near the hill.
  \item \emph{Recall:} say the Marwadi for a kos, then the Marwadi for a traveller, then say \emph{ḍero} again
\end{itemize}

\begin{tcolorbox}[breakable,colback=teal!4,colframe=teal!35!black,title={Before you move on}]
What does \mw{डेरो} mean? (\textquotedblleft{}A camp; a stopping place for the night\textquotedblright{}.) Say it once more.
\end{tcolorbox}
